% GENERATED FILE. Edit canonical lessons, then run npm run generate:books.
% canonical-source-hash: fnv1a64:4be57b5e29eecaf5
% canonical-lessons: GE-C289-ausrutschen, GE-C289-stolpern, GE-C289-stuerzen, GE-C289-sturz, GE-C289-bluten

\chapter{Ein Sturz}
\label{ch:ge-ein-sturz}
\hlchaptermodality{289}

\begin{chapteropening}
\textbf{By the end of this chapter:} \emph{I can say how someone fell and got hurt: to slip, to trip, to fall, a fall and to bleed.}

Complete the last lesson of chapter 289: I can say how someone fell and got hurt: to slip, to trip, to fall, a fall and to bleed.
\end{chapteropening}

\section[ausrutschen]{ausrutschen — to slip}

\label{lesson:GE-C289-ausrutschen}



\begin{quote}
Before the new one: say the German for to have an accident, then the German for the damage.
\end{quote}

\subsection*{You'll want to know: ausrutschen}
\textbf{ausrutschen} — \textquotedblleft{}to slip\textquotedblright{}.

\begin{quote}
Sie rutschte auf dem nassen Boden aus.
\end{quote}

\begin{quote}
\emph{She slipped on the wet floor.}
\end{quote}

Separable, Perfekt with \emph{sein}: \emph{Ich bin ausgerutscht}.

\begin{cousinweb}
\emph{ausrutschen} joins \emph{aus} and \emph{rutschen} \textquotedblleft{}to slide\textquotedblright{}, which is usually traced to Middle High German \emph{rutschen}; EtymDB lists no parent.
\end{cousinweb}

\begin{grammarlens}[title={What you've built}]
One more word for telling, explaining and describing at B1.
\end{grammarlens}

\subsection*{Your turn}
\begin{itemize}
\raggedright
  \item \emph{Say it:} \emph{ausrutschen}
  \item \emph{Say it:} \emph{ausrutschen}, once more
  \item \emph{Say it:} a sentence of your own with \emph{ausrutschen}
  \item \emph{Recall:} say the German for to have an accident, then the German for the damage, then say \emph{ausrutschen} again
\end{itemize}

\begin{tcolorbox}[breakable,colback=teal!4,colframe=teal!35!black,title={Before you move on}]
What does ausrutschen mean? (\textquotedblleft{}To slip\textquotedblright{}.) Say it once more.
\end{tcolorbox}
\section[stolpern]{stolpern — to trip, to stumble}

\label{lesson:GE-C289-stolpern}



\begin{quote}
Before the new one: say the German for the damage, then the German for to slip.
\end{quote}

\subsection*{You'll want to know: stolpern}
\textbf{stolpern} — \textquotedblleft{}to trip, to stumble\textquotedblright{}.

\begin{quote}
Er stolperte über einen Stein.
\end{quote}

\begin{quote}
\emph{He tripped over a stone.}
\end{quote}

Perfekt with \emph{sein}: \emph{Er ist gestolpert}. \emph{Über etwas stolpern}.

\begin{cousinweb}
\emph{stolpern} may be related to Middle High German \emph{stolzen} \textquotedblleft{}to limp\textquotedblright{}, but EtymDB lists no parent and the origin is not settled.
\end{cousinweb}

\begin{grammarlens}[title={What you've built}]
One more word for telling, explaining and describing at B1.
\end{grammarlens}

\subsection*{Your turn}
\begin{itemize}
\raggedright
  \item \emph{Say it:} \emph{stolpern}
  \item \emph{Say it:} \emph{stolpern}, once more
  \item \emph{Say it:} a sentence of your own with \emph{stolpern}
  \item \emph{Recall:} say the German for the damage, then the German for to slip, then say \emph{stolpern} again
\end{itemize}

\begin{tcolorbox}[breakable,colback=teal!4,colframe=teal!35!black,title={Before you move on}]
What does stolpern mean? (\textquotedblleft{}To trip, to stumble\textquotedblright{}.) Say it once more.
\end{tcolorbox}
\section[stürzen]{stürzen — to fall (heavily); to rush}

\label{lesson:GE-C289-stuerzen}



\begin{quote}
Before the new one: say the German for to slip, then the German for to trip.
\end{quote}

\subsection*{You'll want to know: stürzen}
\textbf{stürzen} — \textquotedblleft{}to fall (heavily); to rush\textquotedblright{}.

\begin{quote}
Der alte Mann stürzte auf der Treppe.
\end{quote}

\begin{quote}
\emph{The old man fell on the stairs.}
\end{quote}

Perfekt with \emph{sein}: \emph{Er ist gestürzt}. Also 'to rush': \emph{Sie stürzte aus dem Zimmer}.

\begin{cousinweb}
\emph{stürzen} is inherited from Old High German \emph{sturzen} \textquotedblleft{}to hurl, to plunge, to turn upside down\textquotedblright{} (EtymDB).
\end{cousinweb}

\begin{grammarlens}[title={What you've built}]
One more word for telling, explaining and describing at B1.
\end{grammarlens}

\subsection*{Your turn}
\begin{itemize}
\raggedright
  \item \emph{Say it:} \emph{stürzen}
  \item \emph{Say it:} \emph{stürzen}, once more
  \item \emph{Say it:} a sentence of your own with \emph{stürzen}
  \item \emph{Recall:} say the German for to slip, then the German for to trip, then say \emph{stürzen} again
\end{itemize}

\begin{tcolorbox}[breakable,colback=teal!4,colframe=teal!35!black,title={Before you move on}]
What does stürzen mean? (\textquotedblleft{}To fall (heavily); to rush\textquotedblright{}.) Say it once more.
\end{tcolorbox}
\section[der Sturz]{der Sturz — the fall}

\label{lesson:GE-C289-sturz}



\begin{quote}
Before the new one: say the German for to trip, then the German for to fall.
\end{quote}

\subsection*{You'll want to know: der Sturz}
\textbf{der Sturz} — \textquotedblleft{}the fall\textquotedblright{}.

\begin{quote}
Nach dem Sturz konnte sie nicht mehr gehen.
\end{quote}

\begin{quote}
\emph{After the fall she could not walk any more.}
\end{quote}

Plural: \emph{die Stürze}. Also the fall of a government.

\begin{cousinweb}
\emph{Sturz} is built on the root of \emph{stürzen} \textquotedblleft{}to plunge\textquotedblright{}, Old High German \emph{sturzen} (EtymDB).
\end{cousinweb}

\begin{grammarlens}[title={What you've built}]
One more word for telling, explaining and describing at B1.
\end{grammarlens}

\subsection*{Your turn}
\begin{itemize}
\raggedright
  \item \emph{Say it:} \emph{der Sturz}
  \item \emph{Say it:} \emph{der Sturz}, once more
  \item \emph{Say it:} a sentence of your own with \emph{der Sturz}
  \item \emph{Recall:} say the German for to trip, then the German for to fall, then say \emph{der Sturz} again
\end{itemize}

\begin{tcolorbox}[breakable,colback=teal!4,colframe=teal!35!black,title={Before you move on}]
What does der Sturz mean? (\textquotedblleft{}The fall\textquotedblright{}.) Say it once more.
\end{tcolorbox}
\section[bluten]{bluten — to bleed}

\label{lesson:GE-C289-bluten}



\begin{quote}
Before the new one: say the German for to fall, then the German for the fall.
\end{quote}

\subsection*{You'll want to know: bluten}
\textbf{bluten} — \textquotedblleft{}to bleed\textquotedblright{}.

\begin{quote}
Sein Knie blutete stark.
\end{quote}

\begin{quote}
\emph{His knee bled heavily.}
\end{quote}

Built on \emph{das Blut}. Regular: \emph{blutete}, \emph{hat geblutet}.

\begin{cousinweb}
\emph{bluten} is inherited from Middle High German \emph{bluoten} (EtymDB), built on \emph{Blut} \textquotedblleft{}blood\textquotedblright{}, Proto-Germanic \emph{blōþą}.
\end{cousinweb}

\begin{grammarlens}[title={What you've built}]
One more word for telling, explaining and describing at B1.
\end{grammarlens}

\subsection*{Your turn}
\begin{itemize}
\raggedright
  \item \emph{Say it:} \emph{bluten}
  \item \emph{Say it:} \emph{bluten}, once more
  \item \emph{Say it:} a sentence of your own with \emph{bluten}
  \item \emph{Recall:} say the German for to fall, then the German for the fall, then say \emph{bluten} again
\end{itemize}

\begin{tcolorbox}[breakable,colback=teal!4,colframe=teal!35!black,title={Before you move on}]
What does bluten mean? (\textquotedblleft{}To bleed\textquotedblright{}.) Say it once more.
\end{tcolorbox}
